\documentclass[conference]{IEEEtran}
\usepackage[utf8]{inputenc}
\usepackage[spanish,es-noshorthands]{babel}
\usepackage[T1]{fontenc}
\usepackage{amsmath,amssymb}
\usepackage{graphicx}
\usepackage{booktabs}
\usepackage{listings}
\usepackage{xcolor}
\usepackage{hyperref}
\usepackage{algorithm}
\usepackage{algpseudocode}

\lstset{
  basicstyle=\ttfamily\scriptsize,
  breaklines=true,
  frame=single,
  columns=fullflexible,
  backgroundcolor=\color{gray!8},
  captionpos=b,
}

\begin{document}

\title{Sistema de Rutas y Gesti\'on de Atracciones para un Parque Tem\'atico basado en \'Arboles AVL y Grafos Ponderados}

\author{
  \IEEEauthorblockN{Alfredo Quispe Huamanttupa, Luis Angel Gutierrez Cueva}
  \IEEEauthorblockA{
    Escuela Profesional de Ingenier\'ia Inform\'atica y de Sistemas\\
    Universidad Nacional de San Antonio Abad del Cusco\\
    Curso: Programaci\'on II\\
    Docente: Tany Villalba Villalba
  }
}

\maketitle

\begin{abstract}
Este trabajo presenta el dise\~no e implementaci\'on de un sistema de gesti\'on
para un parque de atracciones cuyo n\'ucleo algor\'itmico se apoya exclusivamente
en dos estructuras de datos implementadas desde cero: un \'arbol AVL gen\'erico
y un grafo ponderado no dirigido. El AVL se utiliza para indexar las
atracciones por nombre y por popularidad, resolviendo b\'usquedas y rankings en
tiempo $O(\log n)$. El grafo modela el mapa del parque y resuelve, mediante el
algoritmo de Dijkstra, dos variantes de ruta \'optima (con y sin tiempo de
espera en fila) y, mediante una b\'usqueda por anchura adaptada a aristas
ponderadas, la consulta de atracciones cercanas dentro de un l\'imite de
distancia. MongoDB se emplea \'unicamente como mecanismo de persistencia: el
servidor construye ambas estructuras en memoria al iniciar y nunca delega en
consultas de base de datos la resoluci\'on de b\'usquedas, rankings o rutas. Se
presentan pruebas unitarias exhaustivas de ambas estructuras y mediciones
emp\'iricas de altura del AVL y tiempo de ejecuci\'on de Dijkstra que confirman
el comportamiento asint\'otico esperado.
\end{abstract}

\begin{IEEEkeywords}
\'Arbol AVL, grafo ponderado, algoritmo de Dijkstra, b\'usqueda en anchura,
estructuras de datos, sistemas de rutas.
\end{IEEEkeywords}

\section{Introducci\'on}

La gesti\'on operativa de un parque de atracciones involucra dos problemas
recurrentes: (1) permitir a los visitantes encontrar r\'apidamente una
atracci\'on por nombre o descubrir las m\'as populares, y (2) calcular c\'omo
desplazarse eficientemente entre atracciones considerando tanto el tiempo de
caminata como el tiempo de espera en fila. Ambos problemas son, en esencia,
problemas cl\'asicos de estructuras de datos: el primero se resuelve de forma
\'optima con un \'arbol de b\'usqueda balanceado, y el segundo con un grafo
ponderado y algoritmos de camino m\'inimo.

El objetivo de este proyecto, desarrollado para el curso de Programaci\'on II,
es demostrar el dominio de estas dos estructuras implement\'andolas
\emph{desde cero} --- sin recurrir a bibliotecas externas de \'arboles, montes
(\emph{heaps}) o grafos--- e integr\'andolas en un sistema completo con
persistencia real (MongoDB), una API REST (Express) y una interfaz web
(React). El requisito de dise\~no no negociable del proyecto es que toda
consulta de negocio (b\'usqueda, ranking, ruta) se resuelva \'unicamente contra
las estructuras en memoria; MongoDB act\'ua como fuente de verdad persistente
pero nunca como motor de consulta.

\section{Marco Te\'orico}

\subsection{\'Arboles AVL}

Un \'arbol AVL \cite{adelson1962} es un \'arbol binario de b\'usqueda (ABB)
autobalanceado en el que, para cada nodo, la diferencia de alturas entre sus
subárboles izquierdo y derecho (\emph{factor de balance}) nunca excede 1 en
valor absoluto:

\begin{equation}
FB(n) = altura(n.izquierdo) - altura(n.derecho), \quad FB(n) \in \{-1,0,1\}
\end{equation}

Esta invariante garantiza que la altura del \'arbol se mantenga en
$O(\log n)$, lo que a su vez acota en $O(\log n)$ el costo de las operaciones
de b\'usqueda, inserci\'on y eliminaci\'on --- a diferencia de un ABB simple,
que degenera a $O(n)$ en el peor caso (por ejemplo, al insertar claves ya
ordenadas).

Cuando una inserci\'on o eliminaci\'on rompe la invariante de balance, se
restaura mediante \emph{rotaciones}, que reorganizan localmente un subárbol
sin alterar el orden in-order de las claves:

\begin{itemize}
  \item \textbf{Rotaci\'on simple derecha (caso Izquierda-Izquierda):} se
    aplica cuando el desbalance proviene del hijo izquierdo del hijo
    izquierdo.
  \item \textbf{Rotaci\'on simple izquierda (caso Derecha-Derecha):} sim\'etrica
    a la anterior.
  \item \textbf{Rotaci\'on doble izquierda-derecha (caso Izquierda-Derecha):}
    una rotaci\'on izquierda sobre el hijo izquierdo seguida de una rotaci\'on
    derecha sobre el nodo.
  \item \textbf{Rotaci\'on doble derecha-izquierda (caso Derecha-Izquierda):}
    sim\'etrica a la anterior.
\end{itemize}

Cada rotaci\'on se ejecuta en tiempo $O(1)$ y, tras una inserci\'on o
eliminaci\'on, se requiere a lo sumo una rotaci\'on (simple o doble) por
nivel del camino de retorno hacia la ra\'iz, preservando la cota $O(\log n)$
total.

\subsection{Grafos ponderados y algoritmo de Dijkstra}

Un grafo ponderado no dirigido $G = (V, E, w)$ se representa en este trabajo
mediante \emph{lista de adyacencia}: cada v\'ertice mantiene un mapa hacia sus
vecinos con el peso de la arista correspondiente. Esta representaci\'on es
adecuada para grafos dispersos (como el mapa de un parque, donde cada
atracci\'on se conecta con un n\'umero peque\~no de atracciones cercanas) y
permite recorrer los vecinos de un nodo en $O(\text{grado}(v))$.

El algoritmo de Dijkstra \cite{dijkstra1959} calcula el camino de costo
m\'inimo desde un v\'ertice origen hacia todos los dem\'as v\'ertices de un
grafo con pesos no negativos, mediante relajaciones sucesivas de las
distancias tentativas. La variante implementada en este proyecto es la
versi\'on cl\'asica $O(V^2)$: en cada una de las $V$ iteraciones se realiza una
b\'usqueda lineal $O(V)$ del v\'ertice no visitado con menor distancia
tentativa, en lugar de emplear una cola de prioridad (\emph{min-heap}), que
reducir\'ia la complejidad a $O((V+E)\log V)$. Para el tama\~no t\'ipico de un
mapa de parque tem\'atico (decenas de nodos), la diferencia pr\'actica entre
ambas variantes es despreciable, mientras que la versi\'on $O(V^2)$ es
sustancialmente m\'as simple de implementar y verificar correctamente sin
bibliotecas externas de montes.

\subsection{B\'usqueda en anchura (BFS) con l\'imite de distancia}

La b\'usqueda en anchura cl\'asica explora un grafo por niveles usando una
cola FIFO y es \'optima \'unicamente en grafos \emph{no ponderados} (donde la
distancia es el n\'umero de saltos). Dado que el mapa del parque es ponderado,
la funcionalidad de ``atracciones cercanas'' se implementa como una variante
de BFS que acumula el peso real de cada arista y relaja la distancia de un
v\'ertice cada vez que se descubre un camino m\'as corto hacia \'el
--- reencolando el v\'ertice cuando esto ocurre, al estilo del algoritmo SPFA.
Esta variante conserva el esp\'iritu de BFS (exploraci\'on dirigida por una
cola, sin estructura de prioridad) y produce distancias acumuladas correctas
sobre aristas ponderadas, descartando toda rama cuya distancia acumulada
exceda el l\'imite solicitado.

\section{Arquitectura del Sistema}

La Figura~\ref{fig:arquitectura} resume el flujo de datos del sistema. Al
iniciar el proceso de Express, se leen desde MongoDB las colecciones
\texttt{Atraccion} y \texttt{Camino} y se construyen tres estructuras en
memoria: un \texttt{AVLTree} indexado por nombre, un \texttt{AVLTree}
indexado por popularidad y un \texttt{GrafoParque}. A partir de ese momento,
\emph{toda} la l\'ogica de negocio expuesta por la API (b\'usqueda, ranking,
c\'alculo de rutas, atracciones cercanas, planificaci\'on de recorridos) se
resuelve exclusivamente contra estas estructuras. MongoDB solo vuelve a
consultarse en operaciones de escritura (\emph{write-through}): cuando se
registra una visita, se actualiza un tiempo de espera, o se agrega una
atracci\'on o camino, el cambio se persiste en MongoDB \emph{y} se refleja en
las estructuras en memoria en la misma operaci\'on.

\begin{figure}[htbp]
\centering
\fbox{\begin{minipage}{0.44\textwidth}
\centering
\vspace{4pt}
\texttt{MongoDB (persistencia)}\\
$\big\downarrow$ \footnotesize carga \'unica al iniciar\\[4pt]
\texttt{AVLTree(nombre) + AVLTree(popularidad) + GrafoParque}\\
\footnotesize (clase \texttt{Parque.js}, en memoria)\\
$\big\downarrow$ \footnotesize toda la l\'ogica pasa por aqu\'i\\[4pt]
\texttt{API REST (Express)}\\
$\big\downarrow$ \footnotesize HTTP/JSON\\[4pt]
\texttt{Frontend (React + Vite + Tailwind)}\\
\vspace{4pt}
\end{minipage}}
\caption{Flujo de datos: MongoDB alimenta las estructuras en memoria una sola
vez al iniciar; la API nunca consulta MongoDB para resolver l\'ogica de
negocio.}
\label{fig:arquitectura}
\end{figure}

Un punto de dise\~no relevante es c\'omo se maneja la actualizaci\'on de la
popularidad de una atracci\'on (al registrar una visita). Dado que la clave
del \texttt{AVLTree} indexado por popularidad cambia, \emph{no} se muta el
nodo in-place --- lo cual violar\'ia la propiedad de orden del \'arbol--- sino
que se elimina el nodo con la clave (popularidad, id) anterior y se reinserta
con la nueva clave. Para evitar colisiones cuando dos atracciones comparten
la misma popularidad, la clave de este \'arbol es siempre el par
\emph{(popularidad, id)}, con desempate lexicogr\'afico por \texttt{id};
esto garantiza que cada clave sea \'unica sin necesidad de listas enlazadas
por nodo.

\subsection{Despliegue en producci\'on}

El sistema se encuentra desplegado y accesible p\'ublicamente, replicando en
la nube la misma arquitectura de tres capas de la Figura~\ref{fig:arquitectura}:

\begin{itemize}
  \item \textbf{Backend (API Express):} desplegado en Render
    (\url{https://render.com}) como \emph{web service} gratuito. Se construye
    con \texttt{npm install} y se ejecuta con \texttt{npm start}; la variable
    de entorno \texttt{MONGODB\_URI} apunta al cl\'uster de MongoDB Atlas.
  \item \textbf{Frontend (React):} desplegado en Vercel
    (\url{https://vercel.com}) como sitio est\'atico, generado con
    \texttt{npm run build}. La variable de entorno \texttt{VITE\_API\_URL},
    fijada en tiempo de \emph{build}, apunta a la URL p\'ublica del backend.
  \item \textbf{Base de datos:} MongoDB Atlas (tier gratuito M0), con la regla
    de red \texttt{0.0.0.0/0} habilitada para aceptar conexiones desde las
    direcciones IP din\'amicas de Render.
\end{itemize}

Cada \texttt{git push} a la rama principal del repositorio dispara un
redespliegue autom\'atico tanto en Render como en Vercel (integraci\'on
continua provista por ambas plataformas sobre el mismo repositorio de
GitHub), sin necesidad de configuraci\'on adicional de CI/CD. Una limitaci\'on
conocida del tier gratuito de Render es que el proceso se suspende tras
\textasciitilde15 minutos de inactividad; la primera solicitud posterior a la
suspensi\'on tarda 30--50 segundos en responder mientras el servicio se
reactiva, lo cual es aceptable para un proyecto acad\'emico pero no para un
entorno de producci\'on real.

\section{Dise\~no e Implementaci\'on}

\subsection{Modelo de datos}

Cada atracci\'on se representa con el esquema
\texttt{\{id, nombre, zona, popularidad, tiempoEsperaFila, coordenadas\}},
donde \texttt{coordenadas} $=\{x,y\}$ fija la posici\'on del nodo en el mapa
SVG del frontend. Cada camino se representa como una arista no dirigida
\texttt{\{origen, destino, tiempoCaminata\}}. El tiempo de espera en fila
vive en el \emph{nodo} (la atracci\'on), no en la arista, reflejando que
esperar en la fila de una atracci\'on es un costo asociado a \emph{llegar a
esa atracci\'on}, independientemente de por d\'onde se llegue.

\subsection{\'Arbol AVL gen\'erico}

La clase \texttt{AVLTree} (Listado~\ref{lst:avl}) es completamente gen\'erica:
recibe un \texttt{comparator(a, b)} en el constructor y no asume nada sobre
el tipo de la clave, lo que permite reutilizar la misma implementaci\'on tanto
para el \'indice por nombre (clave: cadena, orden lexicogr\'afico) como para
el \'indice por popularidad (clave: tupla compuesta). Expone:

\begin{itemize}
  \item \texttt{insert(key, value)}, \texttt{delete(key)}, \texttt{search(key)}
    --- $O(\log n)$ garantizado mediante las cuatro rotaciones descritas en la
    Secci\'on~II-A.
  \item \texttt{inOrder()} --- recorrido in-order completo, usado para el
    listado alfab\'etico de atracciones; $O(n)$.
  \item \texttt{topN(n)} --- recorrido in-order \emph{invertido}
    (derecha-nodo-izquierda) con corte temprano al alcanzar $n$ elementos,
    usado para el ranking de popularidad.
\end{itemize}

\begin{lstlisting}[language=JavaScript, caption={Rotaci\'on simple derecha (caso Izquierda-Izquierda) del AVL.}, label={lst:avl}]
_rotateRight(y) {
  const x = y.left;
  const t2 = x.right;
  x.right = y;
  y.left = t2;
  this._updateHeight(y);
  this._updateHeight(x);
  return x;
}
\end{lstlisting}

\subsection{Grafo del parque y algoritmos de ruta}

La clase \texttt{GrafoParque} implementa:

\begin{itemize}
  \item \texttt{dijkstra(origen, destino, \{incluirEspera\})} --- versi\'on
    $O(V^2)$ descrita en la Secci\'on~II-B, parametrizada por la funci\'on de
    costo. Cuando \texttt{incluirEspera = true} (``ruta \'optima''), el peso
    de relajar la arista $(u,v)$ es
    $peso(u,v) = caminata(u,v) + espera(v)$, es decir, se suma el tiempo de
    espera del nodo de \emph{destino} de cada tramo (nunca el del origen,
    pues no tiene sentido esperar en la atracci\'on de la que se est\'a
    partiendo). Cuando \texttt{incluirEspera = false} (``ruta m\'as corta''),
    $peso(u,v) = caminata(u,v)$.
  \item \texttt{atraccionesCercanas(origen, l\'imite)} --- la variante de BFS
    con relajaci\'on de distancias descrita en la Secci\'on~II-C.
  \item \texttt{planificarRecorrido(origen, paradas)} --- heur\'istica golosa
    del vecino m\'as cercano para el problema de visitar un conjunto de
    paradas: en cada paso elige, mediante Dijkstra, la parada no visitada m\'as
    cercana desde la posici\'on actual. Al tratarse de una variante del
    problema del agente viajero (TSP), que es NP-dif\'icil, esta heur\'istica
    no garantiza el \'optimo global, pero produce un recorrido razonable en
    tiempo polinomial ($O(k^2 \cdot V^2)$ para $k$ paradas).
\end{itemize}

\subsection{API REST}

La Tabla~\ref{tab:endpoints} resume los endpoints expuestos. Cada uno delega
\'integramente en la clase \texttt{Parque.js}, que a su vez opera sobre el
AVL o el grafo seg\'un corresponda; ning\'un controlador ejecuta una consulta
de Mongoose para resolver b\'usqueda, ranking o rutas.

\begin{table}[htbp]
\caption{Endpoints principales de la API}
\label{tab:endpoints}
\centering
\begin{tabular}{@{}p{0.13\textwidth}p{0.16\textwidth}p{0.12\textwidth}@{}}
\toprule
\textbf{M\'etodo/Ruta} & \textbf{Estructura usada} & \textbf{Complejidad} \\
\midrule
\texttt{GET /atracciones/buscar} & AVL(nombre) & $O(\log n)$ / $O(n)$ parcial \\
\texttt{GET /atracciones/ranking} & AVL(popularidad) & $O(k + \log n)$ \\
\texttt{GET /atracciones} & AVL(nombre) & $O(n)$ \\
\texttt{GET /rutas/corta} & Grafo (Dijkstra) & $O(V^2)$ \\
\texttt{GET /rutas/optima} & Grafo (Dijkstra) & $O(V^2)$ \\
\texttt{GET /atracciones/cercanas} & Grafo (BFS ponderado) & $O(V \cdot E)$ peor caso \\
\texttt{POST /rutas/planificar} & Grafo (NN + Dijkstra) & $O(k^2 V^2)$ \\
\texttt{POST /atracciones/visita} & AVL(popularidad) & $O(\log n)$ \\
\bottomrule
\end{tabular}
\end{table}

\subsection{Frontend}

El frontend (React + Vite + Tailwind CSS) consume la API mediante
\texttt{fetch} nativo. El mapa del parque se renderiza como un \texttt{SVG}
con coordenadas fijas asignadas manualmente a cada una de las 16 atracciones
del dataset semilla, evitando deliberadamente un layout autom\'atico
(\emph{force-directed}) para priorizar velocidad de desarrollo y legibilidad.
Cuando se calcula una ruta o un recorrido multi-parada, los nodos y aristas
correspondientes se resaltan visualmente sobre el mapa base.

\section{Resultados de pruebas}

\subsection{Pruebas unitarias}

Se implementaron 19 pruebas unitarias (Node.js \texttt{node:test}) que
verifican, para \texttt{AVLTree}: balance tras 1000 inserciones en orden
aleatorio, correcci\'on de b\'usqueda, actualizaci\'on de valor en claves
duplicadas, eliminaci\'on de hojas / nodos con un hijo / nodos con dos hijos,
vaciado completo, recorridos \texttt{inOrder} y \texttt{topN}, manejo de
claves compuestas, y las cuatro rotaciones (LL, RR, LR, RL); y para
\texttt{GrafoParque}: construcci\'on de aristas no dirigidas, ambas variantes
de costo de Dijkstra, manejo de nodos inalcanzables e inexistentes, l\'imite
de distancia en BFS ponderado, y la heur\'istica de planificaci\'on
multi-parada. Las 19 pruebas pasan (Tabla~\ref{tab:tests}).

\begin{table}[htbp]
\caption{Resumen de pruebas unitarias}
\label{tab:tests}
\centering
\begin{tabular}{@{}lccc@{}}
\toprule
\textbf{Suite} & \textbf{Casos} & \textbf{Aprobados} & \textbf{Fallidos} \\
\midrule
AVLTree & 9 & 9 & 0 \\
GrafoParque & 10 & 10 & 0 \\
\midrule
Total & 19 & 19 & 0 \\
\bottomrule
\end{tabular}
\end{table}

\subsection{Altura del AVL vs. cota te\'orica}

Para verificar emp\'iricamente el comportamiento $O(\log n)$, se insertaron
conjuntos de $n$ claves en orden aleatorio y se midi\'o la altura resultante
del \'arbol, compar\'andola con la cota superior te\'orica
$1.4404 \log_2(n+2) - 0.3277$ \cite{knuth1998}. Los resultados
(Tabla~\ref{tab:altura}) muestran que la altura observada se mantiene siempre
por debajo de la cota te\'orica y crece de forma logar\'itmica, incluso para
$n = 100{,}000$.

\begin{table}[htbp]
\caption{Altura del AVL observada vs. cota te\'orica}
\label{tab:altura}
\centering
\begin{tabular}{@{}rccc@{}}
\toprule
\textbf{$n$} & \textbf{Altura observada} & \textbf{Cota te\'orica} & $\log_2 n$ \\
\midrule
16 & 5 & 5.68 & 4.00 \\
100 & 8 & 9.28 & 6.64 \\
1\,000 & 12 & 14.03 & 9.97 \\
10\,000 & 16 & 18.81 & 13.29 \\
100\,000 & 20 & 23.60 & 16.61 \\
\bottomrule
\end{tabular}
\end{table}

Sobre el dataset semilla real (16 atracciones), el \'arbol indexado por
nombre alcanz\'o una altura de 5, consistente con la fila $n=16$ de la
Tabla~\ref{tab:altura}.

\subsection{Tiempos de ejecuci\'on de Dijkstra}

Se midi\'o el tiempo de una consulta \texttt{rutaOptima} sobre el grafo
completo del dataset semilla (16 nodos, 22 aristas) promediando las 240
consultas posibles entre todos los pares origen-destino: el tiempo promedio
por consulta fue de \textbf{0.0417 ms}. Adicionalmente, se gener\'o una
familia de grafos aleatorios dispersos (grado promedio 4) de tama\~no
creciente para observar la tendencia de crecimiento $O(V^2)$ de la
implementaci\'on (Tabla~\ref{tab:dijkstra}, promedio de 30 consultas por
tama\~no).

\begin{table}[htbp]
\caption{Tiempo promedio de una consulta Dijkstra $O(V^2)$}
\label{tab:dijkstra}
\centering
\begin{tabular}{@{}rrc@{}}
\toprule
\textbf{$V$} & \textbf{$E$ (aprox.)} & \textbf{Tiempo promedio} \\
\midrule
16 & 32 & 0.101 ms \\
50 & 100 & 0.211 ms \\
100 & 200 & 0.814 ms \\
300 & 600 & 3.684 ms \\
500 & 1\,000 & 7.493 ms \\
\bottomrule
\end{tabular}
\end{table}

El crecimiento observado es consistente con $O(V^2)$: al multiplicar $V$ por
$\sim\!10$ (de 50 a 500), el tiempo promedio aumenta $\sim\!35\times$, cercano
al factor $100\times$ esperado en el l\'imite asint\'otico puro, con la
diferencia explicada por el overhead constante de la implementaci\'on y el
efecto del \emph{early-exit} al alcanzar el nodo destino. Para el tama\~no
real del parque (16--20 nodos), el tiempo de respuesta es imperceptible para
el usuario (\textless 0.1 ms de c\'omputo puro).

\section{Discusi\'on}

Los resultados confirman que ambas estructuras cumplen su prop\'osito de
dise\~no. El AVL mantiene su altura logar\'itmica incluso a escalas muy
superiores a las necesarias para este dominio (16--20 atracciones), lo que
deja un amplio margen de crecimiento del dataset sin degradaci\'on de
rendimiento. El uso de una clave compuesta \emph{(popularidad, id)} en el
\'indice de ranking evit\'o por completo la necesidad de manejar colisiones o
listas de nodos por clave, a costa \'unicamente de un comparator ligeramente
m\'as complejo.

Respecto a Dijkstra, la elecci\'on de la variante $O(V^2)$ sin cola de
prioridad fue deliberada: para un grafo con decenas de nodos --- como
cualquier mapa realista de un parque tem\'atico --- la diferencia de tiempo
frente a una implementaci\'on con \emph{min-heap} ($O((V+E)\log V)$) es
inferior al milisegundo, mientras que la implementaci\'on sin heap es
sustancialmente m\'as simple de razonar y depurar sin depender de una
estructura de datos adicional (mont\'iculo) implementada desde cero. Esta
decisi\'on dejar\'ia de ser adecuada si el sistema debiera escalar a mapas
con miles de nodos (por ejemplo, un sistema de rutas a nivel de ciudad), caso
en el que la versi\'on con heap ser\'ia necesaria.

Una limitaci\'on reconocida es que la heur\'istica de vecino m\'as cercano
usada en la planificaci\'on multi-parada no garantiza la ruta \'optima global
(el problema subyacente es NP-dif\'icil); se acept\'o esta limitaci\'on por
ser una funcionalidad secundaria (\'ultima en el orden de prioridad del
proyecto) frente al n\'ucleo AVL/Dijkstra.

\section{Conclusiones}

Se dise\~n\'o e implement\'o un sistema completo de gesti\'on de un parque de
atracciones cuya l\'ogica de negocio central --- b\'usqueda, ranking y
c\'alculo de rutas --- se resuelve \'integramente mediante un \'arbol AVL y un
grafo ponderado implementados desde cero, sin bibliotecas externas de
estructuras de datos. Las pruebas unitarias y las mediciones emp\'iricas
confirman el comportamiento asint\'otico te\'orico de ambas estructuras: altura
logar\'itmica del AVL y crecimiento cuadr\'atico controlado de Dijkstra. La
arquitectura de carga \'unica (MongoDB $\to$ memoria) y escritura
\emph{write-through} permite combinar persistencia real con tiempos de
respuesta de consulta pr\'acticamente instant\'aneos.

\section{Trabajo futuro}

\begin{itemize}
  \item \textbf{Dijkstra con cola de prioridad:} reemplazar la b\'usqueda
    lineal $O(V)$ del m\'inimo por un mont\'iculo binario (\emph{min-heap})
    implementado desde cero, reduciendo la complejidad a
    $O((V+E)\log V)$ para soportar mapas de mayor escala.
  \item \textbf{Layout autom\'atico del mapa:} sustituir las coordenadas fijas
    del SVG por un algoritmo de disposici\'on \emph{force-directed}, \'util si
    el n\'umero de atracciones crece de forma din\'amica.
  \item \textbf{Autenticaci\'on por roles:} restringir el panel
    administrativo (alta de atracciones/caminos, actualizaci\'on de tiempos
    de espera) a usuarios autenticados con rol de administrador.
  \item \textbf{Heur\'istica multi-parada mejorada:} explorar variantes como
    2-opt sobre la soluci\'on golosa inicial para acercar el recorrido
    multi-parada al \'optimo del TSP.
\end{itemize}

\begin{thebibliography}{00}
\bibitem{adelson1962} G.~M. Adelson-Velsky and E.~M. Landis, ``An algorithm
  for the organization of information,'' \emph{Proceedings of the USSR
  Academy of Sciences}, vol.~146, pp.~263--266, 1962.
\bibitem{dijkstra1959} E.~W. Dijkstra, ``A note on two problems in connexion
  with graphs,'' \emph{Numerische Mathematik}, vol.~1, no.~1,
  pp.~269--271, 1959.
\bibitem{knuth1998} D.~E. Knuth, \emph{The Art of Computer Programming,
  Volume 3: Sorting and Searching}, 2nd ed. Addison-Wesley, 1998.
\bibitem{cormen2009} T.~H. Cormen, C.~E. Leiserson, R.~L. Rivest, and
  C.~Stein, \emph{Introduction to Algorithms}, 3rd ed. MIT Press, 2009.
\end{thebibliography}

\end{document}
